\section{Protein interaction across scales}
\label{sec:book-protein-scales}

The five edges in Chapter~\ref{chap:book-one-edge} arise at different resolutions of the same
protein system.  A residue identity belongs to sequence space, an atomic contact belongs to a
conformational description, a family coupling belongs to an evolutionary sample, a model response
belongs to a fitted computation, and epistasis belongs to an intervention with a measured
endpoint.  Their shared residue indices make comparison possible.  Moving between them requires a
declared coarse graining or estimator.

Let $x\in[q]^L$ be a sequence and let $r$ collect atomic coordinates, solvent, protonation, and any
other physical variables retained by a molecular model.  A microscopic energy first defines a
conditional conformational partition function and free energy,
\begin{equation}
Z_{\mathrm{conf}}(x)
=\int e^{-\beta E_{\mathrm{atom}}(x,r)}\,d\mu_x(r),
\qquad
F_{\mathrm{conf}}(x)
=-\beta^{-1}\log Z_{\mathrm{conf}}(x).
\label{eq:book-scale-conformational-free-energy}
\end{equation}
The integral removes physical degrees of freedom.  Its value depends on the ensemble, measure,
temperature, chemical state, and environmental conditions encoded by $(\mu_x,\beta)$.  A single
predicted structure is one representative coordinate configuration; it does not by itself specify
the conformational integral in Eq.~\ref{eq:book-scale-conformational-free-energy}.

Function first enters through a typed measurement.  Let $e$ index an assay protocol, let $c$ collect
the environment and sequence background held fixed by that protocol, and let $\mathsf P_{x,c}$ denote
the molecular or cellular state law presented to the measurement.  The experimental endpoint is
\begin{equation}
Y_e(x,c)=\mathcal A_e(\mathsf P_{x,c},c)\in\mathcal Y_e.
\label{eq:book-scale-assay-map}
\end{equation}
The codomain $\mathcal Y_e$ may contain a binding constant, catalytic rate, abundance, fluorescence,
growth-curve statistic, or a vector of jointly measured quantities.  Assay identity, units, time
window, and preprocessing are part of the definition.  Structure can mediate $Y_e$ through
$\mathsf P_{x,c}$; expression, localization, kinetics, partners, and cellular context can change the
endpoint while a representative folded geometry remains similar.

Population selection requires a further genotype-to-demography map.  The map may use measured traits,
latent ecological variables, or a full population state; its output has its own units and dynamics.
The main text leaves that map unspecified and uses $P_{\mathrm{evo}}$ for the resulting evolutionary
law.  The Malthusian-rate construction is given in the optional population-selection appendix.

Evolution acts on populations and histories rather than on one isolated structure.  A mutation
process, population model, and genealogy together induce a law on the observed leaves
of an alignment,
\begin{equation}
(E_{\mathrm{atom}},\mu,\{\mathcal A_e\},\text{population},T)
\longmapsto
\mathbb P_{\mathrm{evo}}(X_{\mathcal L}),
\qquad X_{\mathcal L}\in[q]^{R\times L}.
\label{eq:book-scale-evolutionary-law}
\end{equation}
Chapter~\ref{chap:book-phylogeny} develops the tree component of this construction.  Selection,
drift, mutation, recombination, sampling, and shared ancestry can all shape the alignment.  A
sequence-level marginal
$P_{\mathrm{evo}}^{\mathrm{seq}}$ consequently combines several mechanisms.  Its negative-log
coordinate
\begin{equation}
H_{\mathrm{seq}}(x)=-\log P_{\mathrm{evo}}^{\mathrm{seq}}(x)
\label{eq:book-scale-sequence-surprisal}
\end{equation}
is a statistical sequence score.  Connecting it quantitatively to
$F_{\mathrm{conf}}(x)$ requires the intervening assay, demographic, and evolutionary maps to be
specified.

Population selection is a downstream model choice.  The assay endpoint $Y_e$ and the induced sequence
law $P_{\mathrm{evo}}$ can be used without assigning either one a universal fitness coordinate.  A
specific mutation--selection operator, its equilibrium eigenproblem, and reversible Gibbs-like special
cases are collected in Appendix~\ref{app:book-population-selection}.  The main text only carries the
boundary condition: interpreting an assay value or a model score as population fitness requires an
explicit genotype-to-demography map and its units.

The complete path used by this book is therefore
\begin{equation}
\boxed{
\begin{gathered}
\text{atomic law}
\longrightarrow\text{conformational ensemble}
\longrightarrow\text{structure, kinetics, and function}\\
\longrightarrow\text{population and phylogenetic law}
\longrightarrow\text{observed MSA}
\longrightarrow\text{estimated sequence operator}\\
\longrightarrow\text{declared compression or spectrum}
\longrightarrow\text{matched experiment}.
\end{gathered}}
\label{eq:book-protein-scale-chain}
\end{equation}
Each arrow changes the state space or the available information.  A structural contact can explain
part of an evolutionary constraint; a distal functional relation can also survive without a short
geometric edge.  The operator framework begins near the middle of this chain, where sequence data
or a trained model are converted into a typed categorical relation.  Physical and functional
claims are recovered only by supplying the missing upstream derivation or downstream endpoint.

The word \emph{entropy} also changes type along this path.  Four recurring quantities are:
\begin{center}
\small
\begin{tabular}{@{}>{\raggedright\arraybackslash}p{0.25\linewidth}
                >{\raggedright\arraybackslash}p{0.29\linewidth}
                >{\raggedright\arraybackslash}p{0.36\linewidth}@{}}
\toprule
Quantity & Distribution or state space & Question answered \\
\midrule
Conformational entropy & molecular configurations conditional on sequence and environment & how physical probability mass is distributed across conformations \\
Sequence entropy & $P_{\mathrm{evo}}^{\mathrm{seq}}$ on $[q]^L$ & how dispersed the induced sequence ensemble is \\
Phylogenetic path or leaf uncertainty & ancestral histories or $X_{\mathcal L}$ on a tree & how much evolutionary history or leaf state remains uncertain \\
Predictive entropy & one model output simplex under a declared context & how uncertain the fitted predictor is about its normalized variable \\
\bottomrule
\end{tabular}
\end{center}
These quantities may influence one another through Eq.~\ref{eq:book-protein-scale-chain}.  Their
numerical equality would require an additional theorem tying the corresponding probability laws
together.

There is likewise no universal monotonicity law for entropy along evolution.  Mutation can broaden
a sequence distribution, while purifying selection can concentrate it; drift, bottlenecks,
recombination, migration, environmental change, and biased sampling alter the result in different
directions.  A claim of entropy production requires a declared trajectory law, time direction,
boundary conditions, and measure of irreversibility.  The entropy of an observed MSA is a property
of its induced and sampled distribution, not a clock that must increase along every lineage.

A cross-scale claim can now be read through six proof obligations:
\begin{equation}
\boxed{
\begin{gathered}
\text{state space}
\to\text{law}
\to\text{observable}\\
{}\to\text{reference or small parameter}
\to\text{controlled remainder}
\to\text{validated effective claim}.
\end{gathered}}
\label{eq:book-cross-scale-proof-obligations}
\end{equation}
Later chapters use this ledger for phylogenetic nulls, sequence Hamiltonians, spectral modes, and
biological interventions.  In each case the effective object is licensed only to the level reached
by its declared reference, controlled approximation, and matched validation endpoint.

Part II begins at the estimation end of this chain: with one categorical pair table and the
contrasts it can identify.
